% !TeX root = ../applied-mathematics1-companion.tex
% Illustrative attracting affine map; the chapter explains general maps.
\begin{center}
\begin{tikzpicture}[scale=3.7,>=Stealth]
\draw[->] (0,0)--(1.2,0) node[right] {$x$};
\draw[->] (0,0)--(0,1.2) node[above] {$y$};
\draw[muted] (0,0)--(1.1,1.1) node[right] {$y=x$};
\draw[thick,ink] (0,.4)--(1.1,.95) node[right] {$y=f(x)=\tfrac12x+\tfrac25$};
\draw[thick,accent,->] (.1,0)--(.1,.45);
\draw[thick,accent,->] (.1,.45)--(.45,.45);
\draw[thick,accent,->] (.45,.45)--(.45,.625);
\draw[thick,accent,->] (.45,.625)--(.625,.625);
\draw[thick,accent] (.625,.625)--(.625,.7125)--(.7125,.7125)
 --(.7125,.75625)--(.75625,.75625)--(.75625,.778125)--(.778125,.778125);
\fill[ink] (.8,.8) circle (.012) node[above left] {$x^*=4/5$};
\node[below] at (.1,0) {$x_0$};
\end{tikzpicture}

{\small Cobweb steps: vertical to the map, horizontal to the diagonal.}
\end{center}
